% SHARD-ID: AQM-91-HAAH-LAURENT-TORUS-QUANT2
% SHARD-TITLE: Haah Laurent Modules and Torus Quant2
% SHARD-SUMMARY: Compares Haah's Laurent-polynomial module formalism with the algebraic-torus Quant2 shard.
% SHARD-SUMMARY: Identifies the finite periodic bridge where the AQM-90 Cayley square becomes Haah's toric-code matrix.
% SHARD-SUMMARY: Records lessons for expressing arithmetic-field dynamics as Weyl-Heisenberg label constraints.
% SHARD-KEYWORDS: Haah, Laurent polynomials, algebraic torus, Quant2, toric code, Weyl-Heisenberg, finite support

\section{Haah Laurent Modules and Torus Quant2}
\label{sec:haah-laurent-torus-quant2}

\subsection{Status and Source Anchors}

This shard is a comparison shard, not a new theorem about all arithmetic
quantum fields.  The local Haah source manifest
\path{references/lattice_stabilizer_codes/SOURCES.md}, lines 17--29, records
the free-module source and exact TeX locators.  Lines 31--43 record the thesis
source for the toric-code matrix, and lines 87--99 record the two-dimensional
classification source.  The algebraic-torus side is AQM-90 and conventions
\texttt{(ay)}--\texttt{(az)}.

\subsection{Lamport Dependency Skeleton}
\[
\begin{array}{rcl}
D1&:&A_T=k[x^{\pm1},y^{\pm1}],\quad T=\Spec A_T.\\
D2&:&R_{\rm tr}=\mathbb F_p[X^{\pm1},Y^{\pm1}]=\mathbb F_p[\mathbb Z^2].\\
D3&:&x,y\text{ are torus coordinates; }X,Y\text{ are translation operators.}\\
D4&:&G=\langle\alpha\rangle\times\langle\beta\rangle\subset T(k).\\
D5&:&R_G=\mathbb F_p[X^{\pm1},Y^{\pm1}]/(X^{N_x}-1,Y^{N_y}-1).\\
D6&:&P_G=R_G^4\cong
\mathbb F_2^{G\times\{x,y\}}{}_q\oplus
\mathbb F_2^{G\times\{x,y\}}{}_p.\\
L1&:&\sigma_\square\text{ reproduces the AQM-90 square CSS subgroup.}\\
L2&:&\sigma_\square^\dagger\lambda\sigma_\square=0
\text{ is }\partial_1\partial_2=0.\\
L3&:&d:A_T\to\Omega^1_{A_T/k}\text{ is not Haah's excitation map.}\\
T1&:&\text{Use finite Laurent matrices as Weyl-label selection principles.}
\end{array}
\]

\subsection{The Same Laurent Ring in Two Different Roles}

Haah's starting point is translation invariance.  A lattice with translation
group \(\mathbb Z^2\) gives the group algebra
\[
  R_{\rm tr}=\mathbb F_p[\mathbb Z^2]
  =\mathbb F_p[X^{\pm1},Y^{\pm1}].
\]
The exponents record lattice positions.  A translation-invariant finite list
of local Pauli or Weyl data is therefore a finite Laurent matrix.  Haah records
this as a map between free modules over the translation group algebra; the
Pauli module is free over that ring, and commutation is encoded by the
Laurent symplectic form.

AQM-90 starts with the split torus
\[
  A_T=k[x^{\pm1},y^{\pm1}],\qquad T=\Spec A_T.
\]
Here \(x,y\) are invertible coordinate functions, not shifts.  Raw Quant2 is a
closed-point residue Weyl field
\[
  E_T(U)=\bigoplus_{z\in U\cap |T|_{\rm cl}}\kappa(z)^2
\]
with finite support.  Thus the formal algebras are isomorphic as rings only
after forgetting their roles.  The comparison convention is deliberately
notational: lower-case \(x,y\) are coordinates; upper-case \(X,Y\) are
translation operators.

\subsection{The Finite Periodic Bridge}

Choose \(\alpha,\beta\in k^\times\), with orders \(N_x,N_y\), and set
\[
  G=\langle\alpha\rangle\times\langle\beta\rangle\subset T(k).
\]
AQM-90 uses this \(G\) as a finite vertex set.  Haah's periodic lattice with
the same periods uses
\[
  R_G=
  \mathbb F_p[X^{\pm1},Y^{\pm1}]/(X^{N_x}-1,Y^{N_y}-1).
\]
As an \(\mathbb F_p\)-vector space, \(R_G\) has the basis
\(\{X^iY^j:0\le i<N_x,\ 0\le j<N_y\}\), hence one basis vector for each site of
the chosen finite torus subgroup.  Multiplication by \(X\) and \(Y\) is the
shift by \(\alpha\) and \(\beta\).

This is the exact bridge.  It is not the claim that \(R_{\rm tr}\) is the
coordinate ring of \(\Spec A_T\).  It is the finite periodic claim that the
translation quotient \(R_G\) acts on the same site-indexed label module that
AQM-90 built from \(G\).

\subsection{Recovering Haah's Toric Matrix from AQM-90}

Use qubit edge labels and order
\[
  P_G=R_G^4=(q_x,q_y,p_x,p_y).
\]
This is the same label space as
\[
  E_\square(G)=
  \mathbb F_2^{G\times\{x,y\}}{}_q
  \oplus
  \mathbb F_2^{G\times\{x,y\}}{}_p
\]
from AQM-90.  With the edge \(e_x(g)\) running from \(g\) to \(Xg\), the
vertex coboundary at \(g\) touches \(e_x(g)\) and \(e_x(X^{-1}g)\), and also
\(e_y(g)\) and \(e_y(Y^{-1}g)\).  Therefore its Laurent column is
\[
  \begin{pmatrix}
  1+\bar X\\
  1+\bar Y\\
  0\\
  0
  \end{pmatrix}.
\]
The face boundary at \(g\) contains
\[
  e_x(g)+e_y(Xg)+e_x(Yg)+e_y(g),
\]
so its phase-label column is
\[
  \begin{pmatrix}
  0\\
  0\\
  1+Y\\
  1+X
  \end{pmatrix}.
\]
Thus the AQM-90 square complex gives the Haah toric matrix
\[
  \sigma_\square=
  \begin{pmatrix}
  1+\bar X & 0\\
  1+\bar Y & 0\\
  \hline
  0 & 1+Y\\
  0 & 1+X
  \end{pmatrix}.
\]
This is Haah's displayed two-dimensional toric-code matrix, modulo the
orientation and row-order convention fixed in \texttt{(az)}.

\subsection{Weyl-Heisenberg Selections}

Let
\[
  \pi:\operatorname{WH}(P_G)\to P_G
\]
be the label map.  Haah's generating matrix selects the subgroup
\[
  L_{\rm gen}(G)=\operatorname{im}\sigma_\square\subset P_G,
\]
so the selected generated Weyl-Heisenberg elements are
\[
  \operatorname{Sel}_{\rm gen}(G)=\pi^{-1}(L_{\rm gen}(G)).
\]
This is exactly AQM-90's CSS subgroup
\[
  L_\square(G)=
  (\operatorname{im}\partial_1^T\oplus0)
  +(0\oplus\operatorname{im}\partial_2).
\]

Haah also gives a stronger test.  Let
\[
  \epsilon_\square=\sigma_\square^\dagger\lambda_2:P_G\to R_G^2.
\]
Then the labels commuting with all selected generators are
\[
  C_{\rm comm}(G)=\ker\epsilon_\square,
\]
and the corresponding accepted and rejected Weyl-Heisenberg elements are
\[
  \operatorname{Sel}_{\rm comm}(G)=\pi^{-1}(C_{\rm comm}(G)),\qquad
  \operatorname{Rej}_{\rm comm}(G)=\pi^{-1}(P_G\setminus C_{\rm comm}(G)).
\]
The matrix equation
\[
  \sigma_\square^\dagger\lambda_2\sigma_\square=0
\]
is the Haah form of the AQM-90 identity \(\partial_1\partial_2=0\).  The
strong toric exactness statement is
\[
  \ker\epsilon_\square=\operatorname{im}\sigma_\square,
\]
with the usual finite-periodic caveat that finite quotients can create global
logical labels.  For AQM purposes, this exactness equation is a sharper
selection principle than bare isotropy.

\subsection{What Is Not Shared}

The Kähler derivation in AQM-90 is
\[
  d:A_T\to\Omega^1_{A_T/k},\qquad
  d(x^ay^b)=x^ay^b(\bar a\,d\log x+\bar b\,d\log y).
\]
It is a coordinate-ring derivation.  Haah's excitation map is instead
\[
  \epsilon=\sigma^\dagger\lambda,
\]
a finite-difference symplectic map on translation labels.  The factors
\(1+X\), \(1+Y\), \(1+\bar X\), and \(1+\bar Y\) are boundary differences, not
Kähler differentials.  Therefore \(d\) should not be advertised as directly
producing the toric-code boundary operator.

There is still a disciplined comparison to pursue: restrict coordinate
families from \(A_T\) to finite supports \(G\), separately build translation
matrices over \(R_G\), and ask when derivative-defined image predicates are
compatible with Haah-style submodules.  AQM-90 already warns that derivative
predicates on finite residue profiles need fibre saturation.  Haah's method
adds the stronger requirement that stable constraints should preferably be
finitely generated \(R_G\)-submodules.

\subsection{Lessons for Arithmetic Quantum Field Dynamics}

The first lesson is grammatical.  A coarse dynamical principle should be a
submodule or kernel inside a finite Weyl label module, not an informal family
of functions.  The output must be a selected or rejected subset of
\(\operatorname{WH}(P_G)\).

The second lesson is exactness.  AQM-90 verified the square boundary identity.
Haah teaches us to ask for the pair
\[
  \operatorname{im}\sigma\subseteq\ker(\sigma^\dagger\lambda)
\]
and then to measure the quotient
\[
  \ker(\sigma^\dagger\lambda)/\operatorname{im}\sigma.
\]
In future arithmetic-field shards this quotient is the natural place to record
finite-periodic logical Weyl labels.

The third lesson is separation of roles.  The algebraic torus supplies
arithmetic points, residue fields, units, and logarithmic differentials.  Haah
supplies the Laurent-module constraint grammar for translation-invariant
topological code layers.  The productive bridge is
\[
  (T,\alpha,\beta)
  \leadsto G
  \leadsto R_G
  \leadsto \sigma\subset P_G
  \leadsto \pi^{-1}(\operatorname{im}\sigma)
  \text{ and }
  \pi^{-1}(\ker\sigma^\dagger\lambda).
\]
This is the precise commonality we can use without collapsing the arithmetic
torus into a bare periodic lattice.
